\documentclass[11pt]{article}
\usepackage{amsmath,amsthm,amssymb,amsfonts}
\usepackage[margin=1in]{geometry}
\usepackage[colorlinks=true,linkcolor=blue,citecolor=blue,urlcolor=blue]{hyperref}
\usepackage{fancyvrb}

\title{ERURH: A Conditional, Formalized Route to the Riemann Hypothesis}
\author{Robert Duran}
\date{}

\theoremstyle{plain}
\newtheorem{theorem}{Theorem}[section]
\newtheorem{lemma}[theorem]{Lemma}
\newtheorem{proposition}[theorem]{Proposition}
\newtheorem{corollary}[theorem]{Corollary}

\theoremstyle{definition}
\newtheorem{definition}[theorem]{Definition}
\newtheorem{assumption}[theorem]{Assumption}
\newtheorem{remark}[theorem]{Remark}

\newcommand{\Rea}{\operatorname{Re}}
\newcommand{\Imc}{\operatorname{Im}}
\newcommand{\R}{\mathbb{R}}
\newcommand{\C}{\mathbb{C}}

\begin{document}
\maketitle

\begin{abstract}
We present a precise conditional proof of the Riemann Hypothesis for the
completed zeta-like function $\xi_\alpha$ used in the ERURH framework. The
argument reduces RH to a finite set of auditable numerical certificates and a
small collection of classical analytic assumptions about the Riemann zeta
function and its explicit formula. The logical chain is fully formalized in
Lean (Lean~4.25.0-rc2 + mathlib). This paper states the exact hypotheses,
explains the ERU energy/inertia machinery, and gives a rigorous, end-to-end
proof of the conditional theorem.
\end{abstract}

\section{Introduction}
The ERURH program reformulates the Riemann Hypothesis (RH) in terms of a
finite set of numerically checkable certificates. The strategy is to translate
the classical explicit formula for $\psi(x)$ into an ``energy--inertia''
framework on the logarithmic scale $s = \log x$, and then show that the
absence of supercritical exponential modes (modes with real part $>1/2$)
forces the RH for a completed zeta-like function $\xi_\alpha$. The novel
content is encapsulated in finite certificates; the classical analytic inputs
are isolated as explicit assumptions. The entire logical chain is formalized
in Lean.

This paper is a rigorous, stand-alone description of the conditional proof.
We make no unconditional claims; all analytic inputs are enumerated and treated
as hypotheses. The main theorem mirrors the Lean theorem
\texttt{RH\_from\_ERURH\_conditional}.

\section{Notation and ERU observables}
Let $\xi_\alpha$ denote the completed zeta-like function used in the ERURH
alpha layer. Its analytic properties (functional equation, explicit formula,
and RH equivalences) are listed in Assumption~\ref{assm:classical}.

Let $\psi(x) = \sum_{p^k\le x} \log p$ and $E(x) = \psi(x) - x$. Define
\begin{equation}
  s = \log x, \qquad
  \log R(s) = \frac{E(e^s)}{e^s}, \qquad
  J_{\mathrm{rel}}(s) = \frac{d}{ds}\log R(s).
\end{equation}
The explicit formula for $\psi(x)$ can be written on the $s$-scale in the form
\begin{equation}\label{eq:explicit-form}
  \log R(s) = \sum_{\rho} a_\rho e^{(\rho-1)s} + \cdots, \qquad
  J_{\mathrm{rel}}(s) = \sum_{\rho} \frac{\rho-1}{\rho}\, a_\rho e^{(\rho-1)s} + \cdots,
\end{equation}
where $\rho = \beta + i\gamma$ ranges over non-trivial zeros of $\zeta(s)$ and
$\cdots$ denotes terms from trivial zeros and smooth test factors. We introduce
abstract coefficients $b_\rho$ from the explicit formula. For the Lean
formalization we fix a normalization in which the explicit-formula package
identifies these coefficients with the prefactor
\[
  \mathrm{explicit\_b}_\rho(\rho) := \pi^{-\rho/2}\Gamma\!\left(\frac{\rho}{2}\right).
\]
Any additional holomorphic factor (e.g., a factor $H(s)$ of moderate growth or
the prefactor $(\rho-1)/\rho$) is absorbed into the abstract coefficient
$b_\rho$ and the explicit-formula assumption. This normalization matches the
Lean definition \texttt{explicit\_b\_rho\_expression}.

\subsection*{Notation summary}
\begin{itemize}
  \item $\xi_\alpha$: the completed zeta-like function of the ERURH alpha layer.
  \item $s=\log x$, $\log R(s)=E(e^s)/e^s$, $J_{\mathrm{rel}}(s)=d(\log R)/ds$.
  \item $\rho=\beta+i\gamma$: a non-trivial zero of $\zeta(s)$.
  \item $b_\rho$: explicit-formula coefficients associated to $\rho$.
  \item $C_{\mathrm{env}}$: formal RMS envelope constant.
\end{itemize}
\noindent The symbol $\alpha$ is used both for the ERU mode parameter and for
generic exponents (e.g., the admissible range in Theorem~\ref{thm:B}); these
uses are unrelated and should not be confused with the fixed subscript in
$\xi_\alpha$.

\section{Main theorem (conditional)}
\begin{theorem}[Conditional ERURH master theorem]\label{thm:main}
Assume the classical zeta package (Assumption~\ref{assm:classical}), the ERU
explicit formula (Assumption~\ref{assm:explicit}), spectral decay (Assumption
\ref{assm:hb}), the large-sieve estimate (Assumption~\ref{assm:ls}), the RMS
window hypotheses (Assumption~\ref{assm:a1a2}), and the correctness of the
alpha certificates and numeric coverage (Assumption~\ref{assm:certs}). Then
there are no ERU modes with weight $\alpha>1/2$, the ERU strong inertia bound
holds, and the Riemann Hypothesis for $\xi_\alpha$ follows.
\end{theorem}

\noindent The exact Lean statement additionally takes an explicit
\texttt{AxiomsShimAccepted} package and separates \texttt{CertificatesCorrectAlpha}
from \texttt{NumericCoverageAlpha}; see the machine-generated exports below.

\begin{remark}
A proof is given in Section~\ref{sec:proof-main} using the Plan~B RMS-local
route and the gate-closure implications of the alpha certificates.
This theorem is conditional: the Lean development proves the implication under
the analytic packages and numeric certificates listed above, and acceptance of
the result depends on external verification of those analytic proofs and the
gate-certified data.
\end{remark}

\section{Formal verification in Lean}
The conditional ERURH master theorem is formalized in Lean 4 (mathlib). The
exact statement and axiom footprint below are machine-generated and included
in the arXiv bundle.
These files are produced by running \texttt{lake env lean docs/core/PrintMasterStatement.lean} and
\texttt{lake env lean docs/core/PrintAxioms_Numeric.lean}, redirecting the outputs to
\texttt{arxiv_submission/lean_master_statement.txt} and \texttt{arxiv_submission/lean_master_axioms.txt}.

{\small\VerbatimInput{lean_master_statement.txt}}

\noindent Axiom footprint for the master theorem:

{\small\VerbatimInput{lean_master_axioms.txt}}

\noindent Axioms shim bundle (Lean export):

{\small\VerbatimInput{lean_axioms_shim_bundle.txt}}

\noindent Gap statements (machine-generated):

{\small\VerbatimInput{lean_gap_statements.txt}}

\noindent These gap statements are paper-level analytic obligations. We provide
full proofs for the primary legacy window-free route (A1 and A2) in this paper,
and full proofs of the A/B/C analytic inputs in the companion preprint. The
fixed-window threshold control remains a documented analytic gap, as noted in
Assumption~\ref{assm:threshold-control}.

\noindent Discharge map for \texttt{AxiomsShimAccepted}:

\begin{center}
\begin{tabular}{|l|l|}
\hline
Field(s) & Discharge path \\
\hline
\texttt{alphaInterfacesBase}, \texttt{explicit\_alpha\_via\_stages} & Assumption~\ref{assm:explicit} (explicit formula package). \\
\texttt{xi\_bounds\_alpha}, \texttt{xi\_argument\_alpha} and their hypothesis bridges & Classical/citable implications (Assumptions Ledger). \\
\texttt{rh\_from\_E\_alpha\_of\_hypotheses} & Classical/citable implications (Assumptions Ledger). \\
\texttt{h\_inertia\_to\_E}, \texttt{h\_inertia\_of\_E}, \texttt{h\_RH\_to\_E} & Classical/citable implications (Assumptions Ledger). \\
\hline
\end{tabular}
\end{center}

\noindent Route B bridge assumptions (primary legacy window-free or alternative
fixed-window computational) are listed in Section~\ref{sec:no-supercritical} and
are not part of \texttt{AxiomsShimAccepted}.

\noindent The beta layer is now an explicit hypothesis bundle (no frozen numeric input is present
in \texttt{data/releases/erurh-v2-core/formal}); see the Assumptions Ledger section below.

\section{No supercritical ERU modes}\label{sec:no-supercritical}
\begin{proposition}[No supercritical ERU modes (conditional)]\label{prop:no-supercritical}
Assume the Plan~B RMS-local hypotheses (Assumption~\ref{assm:a1a2}) and the
correctness of the alpha certificates and numeric coverage
(Assumption~\ref{assm:certs}). Then for every $\beta > 1/2$, there is no ERU
exponential mode of weight $\beta$.
\end{proposition}

\begin{proof}
Fix $\beta>1/2$ and suppose for contradiction that
\texttt{ERU\_mode\_beta}$(\beta)$ holds.
By Lemma~\ref{lem:a1a2-rms-violation}, there exists a window $w$ with
$\mathrm{RMS}_{\alpha}(w) > C_{\mathrm{env}}$. By
Lemma~\ref{lem:rms-violation-opens-gate}, the renormalization gate is open.
Assumption~\ref{assm:certs} implies the relevant renormalization gates are
closed, a contradiction. Hence \texttt{ERU\_mode\_beta}$(\beta)$ does not hold.
Since $\beta>1/2$ was arbitrary, no supercritical ERU exponential mode can
exist.
\end{proof}

\subsubsection*{Two Route B bridges (primary legacy window-free vs alternative fixed-window).}
\begin{itemize}
  \item \textbf{Legacy window-free bridge (primary route).} The paper-level lemma
  \texttt{A1\_from\_supercritical} (Lemma~\ref{lem:a1-from-supercritical}) supplies
  the A1 mode lower bound in an abstract RMS context. A2-low/tail are supplied by
  Lemma~\ref{lem:a2-from-abc}.
  \item \textbf{Fixed-window bridge (alternative computational route).} The Lean
  hypothesis \texttt{ERURH.Alpha.ModeThresholdControlOnCtxRealWindowFamily}
  asserts that the supercritical-mode threshold can be chosen within the
  \texttt{ctx\_real} window family. This is a paper-level analytic assumption;
  the gate checks only the numeric side conditions (see
  Assumption~\ref{assm:threshold-control}).
\end{itemize}

\begin{center}
\begin{tabular}{llll}
\hline
Route & Lean gap & Discharge & Reference \\
\hline
Primary (legacy window-free) & \texttt{A1\_from\_supercritical} and A2 & Paper & Lemmas~\ref{lem:a1-from-supercritical},~\ref{lem:a2-from-abc} \\
Alternative (fixed-window) & \texttt{ModeThresholdControlOnCtxRealWindowFamily} & Paper + gate & Assm.~\ref{assm:threshold-control} \\
\hline
\end{tabular}
\end{center}

\noindent The main narrative below follows the primary legacy window-free route;
the fixed-window computational bridge is included as a documented alternative.

\begin{lemma}[Exponential dominates polynomial]\label{lem:exp-dominates}
Let $a>0$, $k\ge 0$, and $C>0$. Then there exists $S_1$ such that for all
$S\ge S_1$,
\[
\frac{\exp(aS)}{S^k} > C.
\]
\end{lemma}
\begin{proof}
Set $f(S)=\exp(aS)/S^k$ for $S>0$. Then
$\log f(S)=aS-k\log S\to\infty$ as $S\to\infty$.
Choose $S_1$ so that $\log f(S)> \log C$ for all $S\ge S_1$, hence $f(S)>C$.
\end{proof}

\subsubsection*{Legacy window-free A1 lemma (primary route).}
\begin{lemma}[A1 from supercritical mode (legacy window-free)]
\label{lem:a1-from-supercritical}
Let $\beta\in\mathbb{R}$ with $\beta>1/2$. If \texttt{ERU\_mode\_beta}$(\beta)$ holds,
then \texttt{A1\_mode} holds for the chosen RMS context. Equivalently, this is the
statement \texttt{ERURH.A1\_from\_supercritical}.
\end{lemma}
\begin{proof}
We use the admissibility (cofinality) of the window family
(Definition~\ref{def:admissible-window}) and the elementary fact that
$\exp(aS)/S^2 \to \infty$ for any $a>0$.

\textbf{Step 1 (mode threshold).} By \texttt{ERU\_mode\_beta}$(\beta)$ there exists
$s_0$ such that for all $s\ge s_0$,
$|\log R_\alpha^{\mathrm{mode}}(s)| \ge \exp((\beta-1/2)s)$.
Here \texttt{ERU\_mode\_beta} is interpreted on the supercritical mode component
$\log R_\alpha^{\mathrm{mode}}$; this is the analytic content of the A1 bridge.

\textbf{Step 2 (exp dominates the envelope).} Let
$C_\ast := c_{\mathrm{low}}+c_{\mathrm{tail}}+C_{\mathrm{env}}$. Since
$\beta>1/2$, Lemma~\ref{lem:exp-dominates} (with $a=\beta-1/2$, $k=2$) yields
an $S_1$ such that for all $S\ge S_1$, $\exp((\beta-1/2)S)/S^2 > C_\ast$.

\textbf{Step 3 (choose a large window).} By admissibility
(Definition~\ref{def:admissible-window}), there exists a window
$w=[S,S+L]$ with $S\ge \max\{s_0,S_1,1\}$.
For any $s\in I_w=[S,S+L]$ we have $s\ge S\ge s_0$, hence
$|\log R_\alpha^{\mathrm{mode}}(s)|\ge \exp((\beta-1/2)s)\ge \exp((\beta-1/2)S)$.

\textbf{Step 4 (pointwise to RMS).} By Lemma~\ref{lem:pointwise-to-rms},
$\mathrm{RMS}_{\mathrm{mode}}(w)\ge \exp((\beta-1/2)S)\ge
\exp((\beta-1/2)S)/S^2$. Set $K:=\exp((\beta-1/2)S)/S^2$; then
$K>C_\ast$ by Step 2 and $\mathrm{RMS}_{\mathrm{mode}}(w)\ge K$.
This is exactly \texttt{A1\_mode}.
\end{proof}

\noindent The Buchstab bridge below provides one concrete derivation of
\texttt{ERU\_mode\_beta} and the mode coefficient lower bound; in Lean,
\texttt{A1\_from\_supercritical} is derived from that packaging via
\texttt{A1\_from\_supercritical\_of\_buchstab}.

\subsubsection*{Buchstab bridge (analytic derivation of A1)}
The Buchstab bridge supplies one analytic route to
Lemma~\ref{lem:a1-from-supercritical}. The remaining analytic gap is the bridge
\texttt{ERURH.A1\_from\_supercritical\_buchstab} for a chosen RMS context, together
with the explicit-formula identification \texttt{ERURH.ExplicitBRhoExpression}.
The exact Lean statements are machine-exported below.

{\small\VerbatimInput{lean_gap_statements.txt}}

\paragraph{Definition (Buchstab mode lower bound).}
Fix an RMS context $\mathrm{ctx}$ with window type and $\mathrm{RMS\_mode}$.
For $\beta,\gamma\in\mathbb{R}$, we write
\texttt{ModeRMSLowerBound\_from\_buchstab}$(\mathrm{ctx},\beta,\gamma)$ for the
statement that there exists a structural coefficient $c_{\mathrm{struct}}\neq 0$
and a threshold $S_0$ such that for all $S\ge S_0$ there exists a window $w$ with
\[
\mathrm{RMS\_mode}(w)\ \ge\ 
|c_{\mathrm{struct}}\, m(\beta+i\gamma)|\cdot
\frac{\exp((\beta-1/2)S)}{S^2}.
\]
This is exactly the Lean predicate \texttt{ERURH.ModeRMSLowerBound\_from\_buchstab}.

\begin{lemma}[A1 from supercritical mode (Buchstab bridge)]
\label{lem:a1-from-supercritical-buchstab}
Let $\beta\in\mathbb{R}$ with $\beta>1/2$. If \texttt{ERU\_mode\_beta}$(\beta)$ holds,
then there exists $\gamma\in\mathbb{R}$ such that
\texttt{ModeRMSLowerBound\_from\_buchstab}$(\mathrm{ctx},\beta,\gamma)$ holds.
Equivalently, this is the statement \texttt{ERURH.A1\_from\_supercritical\_buchstab}.
\end{lemma}
\begin{proof}
\textbf{Step 1 (unpack the mode).} By \texttt{ERU\_mode\_beta}$(\beta)$ there exists
$s_0$ such that for all $s\ge s_0$,
$|\log R_\alpha(s)| \ge \exp((\beta-1/2)s)$.

\textbf{Step 2 (explicit formula coefficient; paper assumption).} The explicit
formula package identifies the supercritical mode coefficient with
the explicit prefactor (Assumption~\ref{assm:explicit}); combined with
Lemma~\ref{lem:buchstab-coefficient}, this yields a factorization
$b_\rho=c_{\mathrm{struct}}(\rho)\,m(\rho)$ with $c_{\mathrm{struct}}(\rho)\neq 0$.
This is the central paper-level input in the Buchstab bridge.
The normalization of $b_\rho$ matches the preprint definition after absorbing
holomorphic factors into $H(s)$; see Remark~\ref{rem:Hs-mapping}.

\textbf{Step 3 (non-vanishing of the multiplier).} Lemma~\ref{lem:buchstab-m-nonzero}
gives $m(\beta+i\gamma)\neq 0$ for $\beta>0$, hence the mode coefficient is nonzero.

\textbf{Step 4 (RMS lower bound).} By the definition of \texttt{RMS\_mode} for
the mode contribution and the admissibility of the window family
(Definition~\ref{def:admissible-window}), one can choose windows with
arbitrarily large base $S$ and obtain the stated lower bound with coefficient
$|c_{\mathrm{struct}}\,m(\beta+i\gamma)|$.
This is precisely the predicate
\texttt{ModeRMSLowerBound\_from\_buchstab}$(\mathrm{ctx},\beta,\gamma)$.
\end{proof}

\subsubsection*{Buchstab multiplier non-vanishing (paper bridge for A1).}
Recall the multiplicative averaging operator
\[
(\mathcal{T}h)(u) := \frac{1}{u}\int_{u/3}^{u/2} h(t)\,dt .
\]
For a monomial $h(t)=t^{\rho-1}$ with $\rho=\beta+i\gamma$, one computes
\[
(\mathcal{T}h)(u)=u^{\rho-1}\, m(\rho), \qquad
m(\rho)=\frac{2^{-\rho}-3^{-\rho}}{\rho},
\quad 2^{-\rho}:=\exp(-\rho\log 2).
\]
\begin{lemma}[Buchstab multiplier non-vanishing]\label{lem:buchstab-m-nonzero}
If $\beta>0$ then $m(\beta+i\gamma)\neq 0$.
\end{lemma}
\begin{proof}
Write $2^{-\rho}=\exp(-\rho\log 2)$ and $3^{-\rho}=\exp(-\rho\log 3)$.
Then $|2^{-\rho}|=2^{-\beta}$ and $|3^{-\rho}|=3^{-\beta}$. Since $\beta>0$,
we have $2^{-\beta}\neq 3^{-\beta}$, hence by the reverse triangle inequality
\[
|2^{-\rho}-3^{-\rho}| \ge \big||2^{-\rho}|-|3^{-\rho}|\big|
= |2^{-\beta}-3^{-\beta}|>0.
\]
Thus the numerator is nonzero. The denominator $\rho$ is also nonzero because
$\Re(\rho)=\beta>0$. Therefore $m(\rho)\neq 0$.
\end{proof}

\begin{lemma}[Buchstab coefficient identification]\label{lem:buchstab-coefficient}
  Let $\rho=\beta+i\gamma$ be a nontrivial zero with $\beta>1/2$. Assume the
  explicit-formula identification
  $b_\rho=\mathrm{explicit\_b}_\rho(\rho)$ with
  $\mathrm{explicit\_b}_\rho(\rho)=\pi^{-\rho/2}\Gamma(\rho/2)$ as in the
  explicit formula package. Then the ERU coefficient satisfies
  \[
  b_\rho = c_{\mathrm{struct}}(\rho)\, m(\rho)
  \quad\text{with}\quad c_{\mathrm{struct}}(\rho)\neq 0.
  \]
\end{lemma}
\begin{proof}[Proof (step-by-step)]
  \textbf{Step 1 (explicit formula identification).} By
  Assumption~\ref{assm:explicit},
  $b_\rho=\mathrm{explicit\_b}_\rho(\rho)=\pi^{-\rho/2}\Gamma(\rho/2)$.
  Since $\Gamma$ has no zeros and $\pi^{-\rho/2}\neq 0$, we obtain
  $b_\rho\neq 0$ for $\beta>0$.

\textbf{Step 2 (define the structural coefficient).} By
Lemma~\ref{lem:buchstab-m-nonzero}, $m(\rho)\neq 0$ for $\beta>0$. Define
$c_{\mathrm{struct}}(\rho) := b_\rho / m(\rho)$. Then
$b_\rho = c_{\mathrm{struct}}(\rho)\, m(\rho)$ by construction.

\textbf{Step 3 (non-vanishing).} Since $b_\rho\neq 0$ and $m(\rho)\neq 0$,
we have $c_{\mathrm{struct}}(\rho)\neq 0$.
\end{proof}

  \paragraph{Lean correspondence.}
  The explicit-formula identification is the Lean predicate
  \texttt{ERURH.ExplicitBRhoExpression}. From it, Lean derives
  \texttt{ERURH.BuchstabCoefficientNonzero} via
  \texttt{buchstab\_coefficient\_nonzero\_of\_explicit\_b\_rho\_expr}. The lemma
  above is the paper-level justification for the explicit identification and
  the derived factorisation.

\paragraph{Consequences for A1.}
If the supercritical mode contribution to $\log R_\alpha$ carries a coefficient
that is proportional to $m(\rho)$ times a structural factor $c_{\mathrm{struct}}\neq 0$
(from the explicit formula package), then the mode coefficient is nonzero.
Lemma~\ref{lem:buchstab-coefficient} provides this factorisation, and
Lemma~\ref{lem:buchstab-m-nonzero} ensures the multiplier is nonzero. This
reduces the analytic content of A1 to the explicit-formula identification and
the Buchstab lower-bound bridge; the exponential growth itself is already
captured by the multiplier and the elementary $\exp$-dominates-polynomial lemma
used in the Lean core.

\paragraph{Window semantics and admissibility.}
The Lean predicate \texttt{A1\_mode} is defined over an abstract
\texttt{RMSLocalContext} with a type of windows and an associated
\texttt{RMS\_mode}. In the paper we model a window as an interval
$w=[S,S+L]$ with $1\ll L\ll S^\alpha$ and $I_w=[S,S+L]$. The only analytic
input used in Steps~3 and~4 is \emph{admissibility}, i.e., the ability to
choose a window with arbitrarily large $S$. We build this cofinality into the
definition of an admissible window family (Definition~\ref{def:admissible-window}),
so it is part of the window model rather than an extra analytic assumption.

\begin{remark}[Does not replace A/B/C]
Lemma~\ref{lem:a1-from-supercritical-buchstab} is an internal ERURH analytic bridge.
It does \emph{not} replace the classical assumptions A/B/C (zero-density,
coefficient growth, and large-sieve style estimates), which are still required
to derive the A2-tail bound and to connect the explicit formula to the ERU
framework.
\end{remark}

\paragraph{Definitions (paper semantics).}
Fix a window $w$ with associated interval $I_w \subset \mathbb{R}$ (paper-level
window; in Lean this is represented abstractly by a `Window` in `RMSLocalContext`).
Define the RMS-mode contribution by
\[
\mathrm{RMS\_mode}(w) :=
\left(\frac{1}{|I_w|}\int_{I_w} |\log R_\alpha^{\mathrm{mode}}(s)|^2\, ds\right)^{1/2},
\]
where $\log R_\alpha^{\mathrm{mode}}$ denotes the supercritical mode component of
$\log R_\alpha$. The total RMS of $\log R_\alpha$ is denoted
$\mathrm{RMS}_\alpha(w)$ and is decomposed in Assumption~\ref{assm:a1a2}.
For Route~B we use the unique window of \texttt{ctx\_real} (since $n\_windows=1$),
written \texttt{ctx\_real\_window} in Lean.

\begin{lemma}[Pointwise-to-RMS bridge]\label{lem:pointwise-to-rms}
If $|\log R_\alpha^{\mathrm{mode}}(s)| \ge K$ for all $s \in I_w$ and $|I_w|>0$, then
$\mathrm{RMS\_mode}(w) \ge K$. This follows by squaring the pointwise bound,
integrating over $I_w$, and taking square roots:
\[
\frac{1}{|I_w|}\int_{I_w} |\log R_\alpha^{\mathrm{mode}}(s)|^2\, ds \ge \frac{1}{|I_w|}\int_{I_w} K^2\, ds = K^2.
\]
\end{lemma}

\paragraph{Lean correspondence.}
\begin{center}
\begin{tabular}{ll}
Paper object & Lean name \\
\hline
$\log R_\alpha^{\mathrm{mode}}(s)$ & paper-only mode component of \texttt{ERURH.logR\_alpha} \\
$w$ (window) & \texttt{w : ctx.Window} for an abstract \texttt{RMSLocalContext} \\
$I_w$ & window interval associated to \texttt{w} (paper-level) \\
$\mathrm{RMS\_mode}(w)$ & \texttt{ctx.RMS\_mode w} \\
\end{tabular}
\end{center}
The Lean development does not define $\log R_\alpha^{\mathrm{mode}}$ explicitly;
instead, its effect is summarized by the hypothesis \texttt{A1\_mode}. The A1
lemma above interprets \texttt{ERU\_mode\_beta} as a bound on the mode component.

\paragraph{Gate/Lean alignment.}
In the window-free legacy route, \texttt{A1\_from\_supercritical} is treated as an
analytic lemma for an abstract RMS context. If one later instantiates this context
with gate data (e.g., \texttt{ctx\_real}), then \texttt{RMS\_mode} and the envelope
closure are discharged by the numeric certificates; the window-free core itself
does not depend on any fixed window family.

\paragraph{Semantic compatibility of \texttt{RMS\_mode}.}
For the fixed-window route, we also track a semantic lower bound
\[
\mathrm{rms\_semantic}(w) := \inf_{s\in I_w} |\log R_\alpha(s)|.
\]
In Lean this is the definition \texttt{Alpha.rms\_semantic w}. The gate provides
a certificate (Lean: \texttt{GeneratedRMSModeBridge.ctx\_real\_rms\_mode\_ge\_semantic})
that for each \texttt{ctx\_real} window,
\texttt{ctx\_real.RMS\_mode w} $\ge$ \texttt{rms\_semantic w}. Together with
Lemma~\ref{lem:pointwise-to-rms}, this ensures the numeric \texttt{RMS\_mode}
used by the gate is a valid lower bound for the analytic RMS in the
fixed-window chain. The legacy window-free route does not use this certificate.

\section{Reproducibility / Gate}
The reproducible pipeline regenerates the numeric artifacts and rebuilds the Lean development:
\begin{verbatim}
python tools/make_rational_bounds.py
python tools/make_certificate_values.py
lake build ERURH
python scripts/verify_gate.py --skip-pytests
\end{verbatim}

The reproducibility data release includes all numeric inputs used by the gate
(rational bounds, certificates, and RMS context JSON). These are regenerated by
the scripts above and are distinct from the analytic assumptions listed in
Assumption~\ref{assm:classical} and the lemma \texttt{A1\_from\_supercritical}.

\paragraph{Reproducibility manifest.}
The bundle includes \texttt{reproducibility\_manifest.txt} with tool versions
and SHA256 hashes of the frozen numeric inputs used by the gate.
{\small\VerbatimInput{reproducibility_manifest.txt}}

\paragraph{Trust boundary.}
The formal logic is checked by the Lean kernel and mathlib. The numeric gate
depends on the Python toolchain in \texttt{tools/} (including interval
arithmetic via \texttt{mpmath.iv}) and OS-level arithmetic used for packaging
(\texttt{tar}, PowerShell). These are the only computational dependencies
outside Lean.

\paragraph{Interval-checked numeric bridge (example).}
The gate verifies the fixed-window RMS bridge using interval arithmetic on the
witness points. The script \texttt{tools/check\_rms\_mode\_bridge.py} reads
\texttt{rms\_mode\_bridge.json} and \texttt{rms\_context.json} from the frozen
release data (bundled under \texttt{data/releases/erurh-v2-core/formal/}) and
checks, for each window, that
\[
| \log R_\alpha(s_{\mathrm{wit}}) | \le u_{\mathrm{wit}} \le \mathrm{RMS\_mode}(w),
\qquad s_{\mathrm{wit}} \in [s_{\min}, s_{\max}],
\]
using interval enclosures (\texttt{mpmath.iv}) with adaptive precision
(\texttt{dps} raised up to 200 if needed). This is a fully reproducible check:
\begin{verbatim}
python tools/check_rms_mode_bridge.py
\end{verbatim}
and it is executed by \texttt{scripts/verify\_gate.py}. The log line
``OK: RMS mode bridge certificate matches RMS context data.'' confirms success.

\section{Assumptions Ledger}
\begin{itemize}
  \item \textbf{A/B/C (analytic packages).} The analytic inputs used by Plan~B: spectral decay (A), large-sieve bounds (B), and RMS tail/window controls (C).
  \item \textbf{Classical / citable.} Zeta explicit-formula setup and classical implications (including the ERU inertia $\leftrightarrow$ $E$-bound links and RH $\Rightarrow$ $E$-bounds). The zero-parameterization objects $\beta,\gamma$ and coefficients $b_\rho$ remain part of this classical package.
  \item \textbf{Gate + certificates (alpha + RMS + beta).} The alpha certificates and RMS coverage are discharged by the reproducible pipeline (\texttt{verify\_gate.py}), which regenerates the rational bounds and certificate values and checks the Lean build. The beta inertia bundle is now generated from frozen JSON via \texttt{tools/make\_beta\_certificate.py} and is discharged in the same gate.
\end{itemize}


\subsection*{Classical/citable lemmas used as assumptions}
The Lean bundle \texttt{AxiomsShimAccepted} includes the following classical/citable
implications, recorded as explicit hypotheses:
\begin{itemize}
  \item \texttt{InertiaERU\_alpha\_strong} \Rightarrow \texttt{E\_bound\_strong\_alpha}.
  \item \texttt{E\_bound\_strong\_alpha} \Rightarrow \texttt{InertiaERU\_alpha\_strong}.
  \item \texttt{RiemannHypothesis xiAlpha} \Rightarrow \texttt{E\_bound\_strong\_alpha}.
\end{itemize}

\section{Energy, inertia, and ERU modes}
The ERURH formalism introduces global energy functionals and ``gates'' that
measure whether the observable $J_{\mathrm{rel}}$ (or its smoothed variants)
stays under a prescribed envelope. A key notion is the existence of an
exponential ERU mode.

\begin{definition}[ERU mode]
An \emph{ERU mode of weight $\alpha$} (denoted \texttt{ERU\_mode\_beta $\alpha$}
 in Lean) is an asymptotic component of $\log R(s)$ that grows like
$\exp((\alpha-\tfrac{1}{2})s)$ for large $s$. The ERURH theory is structured so
that the presence of a mode with $\alpha>1/2$ forces energy blow-up beyond the
allowable envelope.
\end{definition}

\begin{definition}[Strong inertia]
\emph{Strong ERU inertia} (Lean: \texttt{InertiaERU\_alpha\_strong}) is a
quantitative upper bound on $J_{\mathrm{rel}}(s)$ of the shape
$|J_{\mathrm{rel}}(s)| \le C e^{-s/2}s^2$ for all sufficiently large $s$.
\end{definition}

A central formal lemma is that the absence of supercritical modes implies
strong inertia, and strong inertia implies RH for $\xi_\alpha$ under the
classical zeta assumptions (Assumption~\ref{assm:classical}).

\section{Certificates and gates}
ERURH isolates the new content into finite certificates, each of which is a
finite list of rational inequalities derived from explicit computations.
The alpha layer uses four certificates:
\begin{itemize}
  \item GlobalEnergyCertificate\_alpha,
  \item KernelBlowupCertificate\_alpha,
  \item BridgeInertiaCertificate\_alpha,
  \item StrongInertiaCertificate\_alpha.
\end{itemize}
Correctness of these certificates is encoded as predicates in Lean. The
certificates ensure that the global energy stays below a fixed envelope and
that any supercritical mode would force a kernel energy blow-up. These
implications are formalized as lemmas that close the ERURH gates.
For the Plan~B route, only the GlobalEnergy and KernelBlowup certificates are
required to close the renormalization gate; the bridge/strong inertia
certificates are used in the direct alpha route.

\section{Classical analytic estimates (Theorems A, B, C)}
This section records the classical analytic estimates for the explicit-formula
coefficients $b_\rho$ that underlie the spectral hypotheses used in ERURH. The
statements match the more detailed companion preprint and are included here to
make the analytic inputs explicit.

\subsection*{Standing hypotheses}
Let $\rho=\beta+i\gamma$ range over the non-trivial zeros of $\zeta(s)$,
counted with multiplicity, and let
\[
  b_\rho := \frac{\rho-1}{\rho}\,G(\rho),
  \qquad
  G(s) := \pi^{-s/2}\Gamma\!\left(\frac{s}{2}\right)H(s).
\]
\noindent\textbf{(H1) (Holomorphic model with gamma factor).}
The auxiliary factor $H(s)$ is holomorphic in a neighbourhood of the closed
strip $0\le \Re s\le 1$, and the coefficients $b_\rho$ are defined by the above
formula (in particular, the completed zeta factor
$\pi^{-s/2}\Gamma(s/2)$ is present in $G(s)$).

\medskip
\noindent\textbf{(H2) (Polynomial vertical growth of $H$).}
There exist constants $C_H>0$ and $A_H\ge 0$ such that
\[
  |H(\sigma+it)| \le C_H(1+|t|)^{A_H}
  \qquad (0\le \sigma\le 1,\ t\in\R).
\]

\medskip
\noindent\textbf{(H3) (Riemann--von Mangoldt and unit-interval zero counts).}
Let $N(T)$ denote the number of non-trivial zeros
$\rho=\beta+i\gamma$ with $0<\beta<1$ and $0<\gamma\le T$, counted with
multiplicity. We use the Riemann--von Mangoldt formula and its consequence
that, for $|y|\ge 3$,
\[
  \#\{\rho:\ y<\Imc(\rho)\le y+1\}\ \ll\ \log(|y|+2),
\]
with an absolute implied constant.

\medskip
\begin{theorem}[Theorem A: Dyadic $L^2$ tail bound]\label{thm:A}
Assume \textbf{(H1)}--\textbf{(H3)}. Then there exist constants $c>0$, $k\ge 0$,
and $T_0\ge 3$ such that, for all $T\ge T_0$,
\[
  \sum_{\substack{\rho \\ T < |\gamma| \le 2T}} |b_\rho|^2
  \;\ll\; T^{2k+1}(\log T)\,e^{-2cT}.
\]
In particular, there exist constants $C>0$ and $A>0$ such that, for all
$T\ge T_0$,
\[
  \sum_{\substack{\rho \\ T < |\gamma| \le 2T}} |b_\rho|^2
  \;\le\; C\,T(\log T)^A.
\]
\end{theorem}

\begin{theorem}[Theorem B: Large-sieve mean-square bound]\label{thm:B}
Assume \textbf{(H3)}. Let $T\ge 3$ and $L\ge 1$, and define
\[
  F(u) = \sum_{\substack{\rho\\ T<|\gamma|\le 2T}} b_\rho\, e^{i\gamma u}.
\]
Then
\[
  \frac{1}{L} \int_0^L |F(u)|^2 \, du
    \;\ll\; (\log T)\sum_{\substack{\rho\\ T<|\gamma|\le 2T}} |b_\rho|^2,
\]
with an absolute implied constant. In particular, this holds in any admissible
range $1\ll L\ll T^\alpha$ with $0<\alpha<1$.
\end{theorem}

\begin{theorem}[Theorem C: Uniform RMS tail bound]\label{thm:C}
Assume \textbf{(H1)}--\textbf{(H3)} and fix $B>0$. Define
\[
  \widetilde F_{\mathrm{tail}}(s)
  \;:=\; s^{-2} \sum_{\substack{|\gamma|>S^B}} b_\rho\, e^{i\gamma s},
  \qquad
  K_{\mathrm{tail}}(S,L)^2
  := \frac{1}{L} \int_{S}^{S+L} \bigl|\widetilde F_{\mathrm{tail}}(s)\bigr|^2\,ds.
\]
Then there exist constants $c'>0$ and $S_0\ge 3$ such that, for all
$S\ge S_0$ and all $L>0$,
\[
  K_{\mathrm{tail}}(S,L)^2 \;\ll\; S^{-4}\exp\!\bigl(-2c' S^{B}\bigr).
\]
In particular, $K_{\mathrm{tail}}(S,L)^2\le M_{\mathrm{tail}}^2$ uniformly in
$S$ and $L$, for a suitable constant $M_{\mathrm{tail}}>0$.
\end{theorem}

\noindent Proofs of Theorems~\ref{thm:A}--\ref{thm:C} are not reproduced here;
see the companion preprint \texttt{Theorem\_ABC\_preprint.tex}. We record only
the statements and use them as analytic inputs.

\begin{lemma}[A2-low and A2-tail from classical inputs]\label{lem:a2-from-abc}
Assume the explicit-formula package (Assumption~\ref{assm:explicit}) together
with Theorems~\ref{thm:A}--\ref{thm:C}. Then there exist constants
$c_{\mathrm{low}}, c_{\mathrm{tail}}$ such that, for every admissible window
$w=[S,S+L]$ with $1\ll L\ll S^\alpha$, one has
\[
  \mathrm{RMS}_{\mathrm{low}}(w) \le c_{\mathrm{low}},
  \qquad
  \mathrm{RMS}_{\mathrm{tail}}(w) \le c_{\mathrm{tail}}.
\]
Equivalently, the A2-low and A2-tail clauses of Assumption~\ref{assm:a1a2}
hold for the chosen RMS context.
\end{lemma}

\begin{proof}
\textbf{A2-low.} Choose a fixed cutoff $T_0$ and decompose the explicit formula
into: (i) finitely many modes with $|\gamma|\le T_0$, (ii) trivial-zero terms,
and (iii) the smooth test-factor contribution. Under
Assumption~\ref{assm:explicit} and the explicit-formula normalization, each of
these pieces is uniformly bounded on admissible windows $[S,S+L]$ with
$1\ll L\ll S^\alpha$. The finite sum therefore yields a uniform constant
$c_{\mathrm{low}}$ for $\mathrm{RMS}_{\mathrm{low}}(w)$. The constant depends
only on the fixed cutoff $T_0$ and the explicit-formula normalization, not on
$S$ or $L$. The cutoff $T_0$ is fixed once in the analytic package (see the
preprint), so $c_{\mathrm{low}}$ is frozen with the choice of kernel and
normalization.

\textbf{A2-tail.} The tail RMS on $w=[S,S+L]$ is, by definition, the normalized
mean square of the tail contribution with the $s^{-2}$ weight. This is exactly
the quantity $K_{\mathrm{tail}}(S,L)$ in Theorem~\ref{thm:C}. Theorem~\ref{thm:C}
is derived from Theorems~\ref{thm:A} and~\ref{thm:B} in the preprint, using the
dyadic $L^2$ tail bound and the large-sieve inequality for separated exponentials.
Thus $K_{\mathrm{tail}}(S,L)^2 \le M_{\mathrm{tail}}^2$ uniformly on admissible
windows, and we take $c_{\mathrm{tail}} := M_{\mathrm{tail}}$.
Theorem~\ref{thm:C} is stated uniformly in $S$ and $L$, so no extra hypotheses
are required beyond Theorems~\ref{thm:A}--\ref{thm:C} and the explicit-formula
package.
No additional hypotheses are used beyond Theorems~\ref{thm:A}--\ref{thm:C} and
the explicit-formula package.
\end{proof}

\noindent Full derivations (including constants and admissible ranges) are
provided in \texttt{Theorem\_ABC\_preprint.tex}, which is included in the arXiv
bundle for this submission.

\paragraph{Lean correspondence for A2.}
The paper inequalities
$\mathrm{RMS}_{\mathrm{low}}(w)\le c_{\mathrm{low}}$ and
$\mathrm{RMS}_{\mathrm{tail}}(w)\le c_{\mathrm{tail}}$ are encoded in Lean as
\texttt{ERURH.A2Low\_RMS ctx} and \texttt{ERURH.A2Tail\_RMS ctx}, where
\texttt{ctx.RMS\_low}, \texttt{ctx.RMS\_tail}, \texttt{ctx.c\_low},
\texttt{ctx.c\_tail} live in \texttt{ERURH.RMSLocalContext}. The decomposition
$\mathrm{RMS}_\alpha \ge \mathrm{RMS}_{\mathrm{mode}}-(\mathrm{RMS}_{\mathrm{low}}+
\mathrm{RMS}_{\mathrm{tail}})$ matches \texttt{ctx.RMS\_decomp}. For the
generated context \texttt{ctx\_real}, the hypotheses are discharged by the
generated lemmas \texttt{Alpha.GeneratedRMSContext.ctx\_real\_A2Low} and
\texttt{Alpha.GeneratedRMSContext.ctx\_real\_A2Tail}.

These theorems provide the analytic content behind the spectral decay
Assumption~\ref{assm:hb} and the large-sieve estimate
Assumption~\ref{assm:ls}. Proofs are given in the companion preprint
\texttt{Theorem\_ABC\_preprint.tex} in the same submission folder, with
references to standard sources (e.g. Titchmarsh, Montgomery--Vaughan, Ivic);
we record only the statements here and treat them as classical inputs.

\section{External analytic assumptions}
The conditional theorem uses the following external packages. We list the
assumptions explicitly to isolate analytic input from ERURH-specific logic.

\begin{assumption}[Classical zeta package]\label{assm:classical}
The completed function $\xi_\alpha$ satisfies the standard analytic properties
of the Riemann zeta function: explicit formula for $\psi(x)$, functional
equation, standard bounds for $\zeta'/\zeta$ in vertical strips, the
Riemann--von Mangoldt zero counting formula (including local counts), and the
classical equivalence between RH and bounds on $E(x)$.
\end{assumption}

\begin{assumption}[Explicit formula for the ERU observable]\label{assm:explicit}
The ERU observable $\log R(s)$ and its derivative $J_{\mathrm{rel}}(s)$ admit an
explicit formula of the form~\eqref{eq:explicit-form} in which the coefficients
$b_\rho$ are identified with the explicit prefactor
$\mathrm{explicit\_b}_\rho(\rho)=\pi^{-\rho/2}\Gamma(\rho/2)$ under the chosen
normalization. Any additional holomorphic factors (e.g., $H(s)$ or
$(\rho-1)/\rho$) are absorbed into $b_\rho$ and the explicit-formula package.
\end{assumption}

\begin{remark}[Mapping to the $H(s)$ factor]\label{rem:Hs-mapping}
In the A/B/C preprint we write
\[
  b_\rho := \frac{\rho-1}{\rho}\,G(\rho),
  \qquad
  G(s) := \pi^{-s/2}\Gamma\!\left(\frac{s}{2}\right)H(s).
\]
For ERURH, the factor $H(s)$ is the Mellin transform of the ERU kernel and
auxiliary weights from the explicit formula. Under the normalization above,
any additional holomorphic factors are absorbed into $b_\rho$ (so the Lean
predicate \texttt{ExplicitBRhoExpression} uses only the gamma prefactor). We
assume $H(\rho)\neq 0$ in the analytic preprint to avoid trivial coefficients,
so the analytic assumptions on $H(s)$ in the preprint (H1–H3) apply verbatim.
\end{remark}

\begin{remark}[Nonvanishing for the Buchstab kernel]\label{rem:buchstab-H-nonzero}
The alpha bridge data use the multiplicative averaging kernel
$(\mathcal{T}h)(u)=\frac{1}{u}\int_{u/3}^{u/2} h(t)\,dt$, whose Mellin multiplier is
$m(s)=(2^{-s}-3^{-s})/s$ (Lemma~\ref{lem:buchstab-m-nonzero}). For $\Re(s)>0$,
Lemma~\ref{lem:buchstab-m-nonzero} gives $m(s)\neq 0$, so the kernel contribution
to $H(s)$ does not vanish on nontrivial zeros. The remaining nonvanishing
requirement in Assumption~\ref{assm:explicit} concerns only auxiliary holomorphic
weights from the analytic preprint; in Lean they are absorbed into $b_\rho$.
\end{remark}

\begin{assumption}[Spectral decay: $H_b$]\label{assm:hb}
The coefficient family $\{b_\rho\}$ satisfies:
(i) a pointwise polynomial (or exponential) bound in $|\gamma|$;
(ii) a dyadic $L^2$ tail bound
\[\sum_{T<|\gamma|\le 2T} |b_\rho|^2 \le C T (\log T)^A\]
for sufficiently large $T$.
\end{assumption}

\begin{assumption}[Large-sieve for zero ordinates]\label{assm:ls}
For each dyadic band $T<|\gamma|\le 2T$ and each $L\ge 1$,
\[\frac{1}{L}\int_0^L\Bigl|\sum_{T<|\gamma|\le 2T} b_\rho e^{i\gamma u}\Bigr|^2\,du
  \ll (\log T)\sum_{T<|\gamma|\le 2T}|b_\rho|^2.\]
This is a standard large-sieve inequality for separated exponentials combined
with local zero counting.
\end{assumption}

\begin{definition}[Admissible window family]\label{def:admissible-window}
A family of windows $w=[S,S+L]$ with $1\ll L\ll S^\alpha$ is called
\emph{admissible} if it is cofinal: for every $S_0$ there exists a window
with $S\ge S_0$ and $1\ll L\ll S^\alpha$.
\end{definition}

\begin{assumption}[RMS window hypotheses]\label{assm:a1a2}
Fix an admissible window family (Definition~\ref{def:admissible-window}) of
windows $w=[S,S+L]$ with $1\ll L\ll S^\alpha$.
There exist RMS quantities $\mathrm{RMS}_{\alpha},\mathrm{RMS}_{\mathrm{mode}},
\mathrm{RMS}_{\mathrm{low}},\mathrm{RMS}_{\mathrm{tail}}$ and constants
$c_{\mathrm{low}},c_{\mathrm{tail}},C_{\mathrm{env}}$ such that:
\begin{itemize}
  \item \textbf{A1 (mode $\Rightarrow$ window bridge):} for every $\beta>1/2$, if
  \texttt{ERU\_mode\_beta}$(\beta)$ holds then there exist an admissible window
  $w$ and a constant $K$ with
  $K>c_{\mathrm{low}}+c_{\mathrm{tail}}+C_{\mathrm{env}}$ such that
  $\mathrm{RMS}_{\mathrm{mode}}(w)\ge K$.
  \item \textbf{A2-low:} for every admissible window $w$,
  $\mathrm{RMS}_{\mathrm{low}}(w)\le c_{\mathrm{low}}$.
  \item \textbf{A2-tail:} for every admissible window $w$,
  $\mathrm{RMS}_{\mathrm{tail}}(w)\le c_{\mathrm{tail}}$.
  \item \textbf{Decomposition:} for every admissible window $w$,
  \[
    \mathrm{RMS}_{\alpha}(w)\ \ge\ \mathrm{RMS}_{\mathrm{mode}}(w)\;-\;
    \mathrm{RMS}_{\mathrm{low}}(w)\;-\;\mathrm{RMS}_{\mathrm{tail}}(w).
  \]
\end{itemize}
These hypotheses are the analytic inputs of the Plan~B RMS-local route. In
particular, A1 is formulated as a \emph{conditional bridge}: it produces a
witness window only under the presence of a putative supercritical mode, and
therefore does not assert an RMS violation in the closed-gate regime.
The admissibility (cofinality) of the window family is part of the window
model, not an extra analytic assumption.
\end{assumption}

\begin{assumption}[Fixed-window threshold control]\label{assm:threshold-control}
Let $S_{\max}$ denote the maximal lower endpoint among the fixed
\texttt{ctx\_real} windows (Lean: \texttt{windowMinMax}). For every
$\beta>1/2$, if \texttt{ERU\_mode\_beta}$(\beta)$ holds then there exists
$s_0\le S_{\max}$ such that for all $s\ge s_0$,
$|\log R_\alpha(s)|\ge \exp((\beta-1/2)s)$. This is precisely the Lean
predicate \texttt{ERURH.Alpha.ModeThresholdControlOnCtxRealWindowFamily}.
\end{assumption}

\begin{remark}
This is a paper-level analytic assumption. The gate certifies the numeric
inequality $S_{\max}\ge s_0$ once the threshold control is supplied, but it does
not prove the analytic implication from \texttt{ERU\_mode\_beta} to that bound.
We provide a discussion of the analytic gap; a full proof is deferred to a
supplement.
\end{remark}

\paragraph{Discussion (analytic gap for Assumption~\ref{assm:threshold-control}).}
\textbf{Step 1.} Unpack \texttt{ERU\_mode\_beta}$(\beta)$ to obtain a threshold
$s_0$ such that $|\log R_\alpha(s)|\ge \exp((\beta-1/2)s)$ for all $s\ge s_0$.

\textbf{Step 2.} Establish an analytic upper bound on $s_0$ in terms of
explicit-formula data (e.g., via a lower bound on the relevant mode coefficient
and its contribution to $\log R_\alpha$). This is the only genuinely analytic
input in the fixed-window bridge.

\textbf{Step 3.} Use the gate-certified numerical inequality
$S_{\max}\ge s_0$ (from \texttt{formal\_report\_analytic.json}) to conclude
$s_0\le S_{\max}$, which is exactly Assumption~\ref{assm:threshold-control}.

\begin{remark}
In the legacy window-free route, A1 is supplied by
Lemma~\ref{lem:a1-from-supercritical-buchstab}, while A2-low and A2-tail are discharged
by Lemma~\ref{lem:a2-from-abc} from the classical inputs.
\end{remark}

\begin{assumption}[Numeric certificates]\label{assm:certs}
The alpha certificates (GlobalEnergy and KernelBlowup) are correct, and the
numeric coverage assumptions required to close the renormalization gates hold.
In particular, the coverage asserts that for every admissible window $w$ in the
covered family (the Lean predicate \texttt{NumericCoverageAlpha}), we have
$\mathrm{RMS}_{\alpha}(w) \le C_{\mathrm{env}}$, i.e., the renormalization gate
is closed on that family.
These are finite rational inequalities generated by the ERURH pipeline.
\end{assumption}

\begin{assumption}[Beta inertia package]\label{assm:beta}
The beta-layer numerical inertia certificates are correct and supply any
additional numeric bounds required by the pipeline. This package is external
and is not used in the core alpha implication, but is included for completeness
of the ERURH pipeline.
\end{assumption}

\section{Plan B: RMS-local route}
The Plan~B proof decomposes the RMS of the ERU observable into mode/low/tail
contributions. The following lemma is formalized in
\texttt{ERURH\_A2Hypotheses.lean}.

\begin{lemma}[Supercritical mode forces an RMS violation]\label{lem:a1a2-rms-violation}
Assume Assumption~\ref{assm:a1a2}. Let $\beta>1/2$. If
\texttt{ERU\_mode\_beta}$(\beta)$ holds, then there exists a window $w$ with
$\mathrm{RMS}_{\alpha}(w) > C_{\mathrm{env}}$.
\end{lemma}

\begin{proof}
Assume $\beta>1/2$ and \texttt{ERU\_mode\_beta}$(\beta)$. By the A1 clause of
Assumption~\ref{assm:a1a2}, choose an admissible window $w$ and $K$ such that
$K>c_{\mathrm{low}}+c_{\mathrm{tail}}+C_{\mathrm{env}}$ and
$\mathrm{RMS}_{\mathrm{mode}}(w)\ge K$. Using the decomposition inequality and
the A2-low/tail bounds from Assumption~\ref{assm:a1a2}, we obtain
\begin{align*}
\mathrm{RMS}_{\alpha}(w)
&\ge \mathrm{RMS}_{\mathrm{mode}}(w)-\mathrm{RMS}_{\mathrm{low}}(w)-\mathrm{RMS}_{\mathrm{tail}}(w) \\
&\ge K - c_{\mathrm{low}} - c_{\mathrm{tail}} \\
&> C_{\mathrm{env}}.
\end{align*}
This is the desired RMS violation.
\end{proof}

A window whose RMS exceeds the envelope forces a renormalization gate to open.
This is a direct consequence of the definition of the gate predicate.

\begin{lemma}[RMS violation opens a gate]\label{lem:rms-violation-opens-gate}
If $\mathrm{RMS}_{\alpha}(w) > C_{\mathrm{env}}$ for some window $w$, then the
renormalization gate is open.
\end{lemma}

The correctness of the certificates (Assumption~\ref{assm:certs}) closes the
renormalization gates; thus the RMS violation is incompatible with the closed
gate regime. Consequently, under Assumptions~\ref{assm:a1a2} and~\ref{assm:certs},
supercritical modes cannot exist.

\section{Proof of the main theorem}\label{sec:proof-main}
We now prove Theorem~\ref{thm:main}. This is the mathematical counterpart of
the Lean theorem \texttt{RH\_from\_ERURH\_conditional}.

\begin{proof}[Proof of Theorem~\ref{thm:main}]
By Proposition~\ref{prop:no-supercritical}, there are no ERU exponential modes
with weight $>1/2$. The Lean lemma
\texttt{InertiaERU\_alpha\_strong\_of\_no\_modes\_via\_certificates} yields
strong ERU inertia from the absence of modes (using
Assumption~\ref{assm:certs}). Finally, the formal equivalence between strong
ERU inertia and RH for $\xi_\alpha$ (under Assumption~\ref{assm:classical})
gives RH.
\end{proof}

\begin{remark}
Assumptions~\ref{assm:explicit}--\ref{assm:ls} are classical analytic inputs
(derivable from explicit formula technology, Stirling estimates, bounds for
$\zeta'/\zeta$, and Riemann--von Mangoldt zero counts). The ERURH proof is a
formal reduction: once these analytic statements and the finite certificates
are established, RH follows mechanically.
\end{remark}

\begin{remark}
In the Lean master statement, the Plan~B window hypotheses are packaged as
\texttt{hA1 : A1\_mode ctx} and \texttt{hLow : A2Low\_RMS ctx}/\texttt{hTail :
A2Tail\_RMS ctx} inside \texttt{WindowScalingAssumptions}. In this paper,
Assumption~\ref{assm:a1a2} is formulated so that the A1 clause is a
\emph{mode-to-window bridge} rather than an unconditional witness. This matches
the intended analytic semantics: A1 produces a window only under a putative
supercritical mode, so it is compatible with the closed-gate regime.
\end{remark}

\begin{remark}
With A1 formulated as an implication from \texttt{ERU\_mode\_beta} to a witness
window, the core Plan~B contradiction has the standard form:
\[
(\exists\,\beta>1/2\ \text{supercritical mode})\ \Rightarrow\
(\exists\,w\ \text{RMS violation})\ \Rightarrow\ \text{gate opens},
\]
which contradicts gate closure from the finite certificates. Thus the
conditional theorem is not vacuous: the hypotheses do not force an RMS
violation unless a supercritical mode exists.
\end{remark}

\section{Formal status and gaps}\label{sec:status-gaps}
The ERURH development separates formal logic from external analytic and numeric
inputs. The current status can be summarized as follows.
\begin{itemize}
  \item \textbf{Formalized in Lean:} definitions of the ERU observables,
  energy/inertia machinery, gate predicates, the Plan~B RMS-local chain, and
  the conditional theorem \texttt{RH\_from\_ERURH\_conditional}.
  \item \textbf{External classical assumptions:} explicit formula, functional
  equation, bounds for $\zeta'/\zeta$, and Riemann--von Mangoldt zero counts
  (Assumption~\ref{assm:classical}), together with the explicit-formula
  framework for $b_\rho$ (Assumption~\ref{assm:explicit}).
  \item \textbf{Analytic inputs:} Theorems~\ref{thm:A}--\ref{thm:C} (spectral
  decay and large-sieve bounds) together with the analytic bridges
  Lemma~\ref{lem:a1-from-supercritical-buchstab} (A1) and Lemma~\ref{lem:a2-from-abc}
  (A2). The remaining gap is the classical derivation of Theorems~\ref{thm:A}--\ref{thm:C}
  from standard zeta estimates. Full proofs of Theorems~\ref{thm:A}--\ref{thm:C}
  are included in the bundled preprint \texttt{Theorem\_ABC\_preprint.tex}.
  \item \textbf{Numeric inputs:} correctness of the finite certificates and
  numeric coverage assumptions (Assumption~\ref{assm:certs}), which are
  generated by the ERURH pipeline as exact rational data.
\end{itemize}

\section{Formalization and reproducibility}
The definitions and implications above are formalized in Lean. The main
logical entry points are:
\begin{itemize}
  \item \texttt{ERURH\_MasterTheoremSummary.lean}:
        \texttt{RH\_from\_ERURH\_conditional}.
  \item \texttt{ERURH\_MasterTheoremPlanB.lean}:
        Plan~B RMS-local route.
  \item \texttt{ERURH\_A2Hypotheses.lean}:
        A1/A2 hypotheses and RMS decomposition.
  \item \texttt{ERURH\_GatesAlpha.lean} and \texttt{ERURH\_GatesFromRMS.lean}:
        gate closure and RMS-to-gate implications.
  \item Reproducibility gate:
        \texttt{scripts/verify\_gate.py} regenerates artifacts and checks the
        Lean build against the current certificate data.
\end{itemize}
The certificates are generated by reproducible scripts and imported as exact
rational data. The formal statement does not rely on floating-point computation.

\section{Conclusion}
ERURH provides a formal, conditional path from explicit-formula analysis to RH.
All non-ERURH analytic content is isolated in a short list of assumptions, and
all ERURH-specific reasoning is formalized. The remaining work is to supply
proofs (or certified bounds) for the analytic assumptions and to validate the
finite certificates with fully verified numerical pipelines.

\appendix
\section{Dependency map}\label{app:dependencies}
This appendix summarizes the logical dependencies among the main components.
Arrows indicate implication or usage.

\subsection*{High-level chain}
\begin{itemize}
  \item Classical zeta inputs (Assumption~\ref{assm:classical}) together with
  the ERU explicit-formula assumption (Assumption~\ref{assm:explicit}) supply
  the explicit-formula layer for $\log R$ and $J_{\mathrm{rel}}$.
  \item Classical estimates (Theorems~\ref{thm:A}--\ref{thm:C})
  $\Rightarrow$ spectral decay and large-sieve inputs
  (Assumptions~\ref{assm:hb}--\ref{assm:ls}).
  \item A1/A2 window hypotheses (Assumption~\ref{assm:a1a2}) plus a
  putative supercritical mode $\texttt{ERU\_mode\_beta}(\beta)$
  $\Rightarrow$ RMS violation $\Rightarrow$ gate opens.
  \item Certificates (Assumption~\ref{assm:certs})
  $\Rightarrow$ gate closed.
  \item RMS violation and gate closed $\Rightarrow$ contradiction; in the
  Plan~B chain this excludes supercritical modes under A1/A2.
  \item No supercritical mode $\Rightarrow$ strong inertia
  $\Rightarrow$ RH for $\xi_\alpha$ (using Assumption~\ref{assm:classical}).
\end{itemize}

\subsection*{Lean file map}
\begin{itemize}
  \item Main theorem:
  \texttt{lean/formal\_proofs/ERURH/ERURH\_MasterTheoremSummary.lean}
  (theorem \texttt{RH\_from\_ERURH\_conditional}).
  \item Plan B chain:
  \texttt{lean/formal\_proofs/ERURH/ERURH\_MasterTheoremPlanB.lean},
  \texttt{lean/formal\_proofs/ERURH/ERURH\_A2Hypotheses.lean}.
  \item Gates:
  \texttt{lean/formal\_proofs/ERURH/ERURH\_GatesAlpha.lean},
  \texttt{lean/formal\_proofs/ERURH/ERURH\_GatesFromRMS.lean}.
  \item Assumption bundles:
  \texttt{lean/formal\_proofs/ERURH/ERURH\_ClassicalZetaAssumptions.lean},
  \texttt{lean/formal\_proofs/ERURH/ERURH\_ExplicitFormulaAlpha.lean},
  \texttt{lean/formal\_proofs/ERURH/ERURH\_LargeSieveGammaSketch.lean}.
  \item Certificates:
  \texttt{lean/formal\_proofs/ERURH/EnergyCertificatesAlpha.lean},
  \texttt{lean/formal\_proofs/ERURH/InertiaCertificatesAlpha.lean}.
\end{itemize}

\section{Proof obligations}\label{app:proof-obligations}
This appendix records the remaining Lean-exported hypotheses required for the
Route~B closure. The machine-generated statements are included below.

{\small\VerbatimInput{lean_gap_statements.txt}}

\noindent A field-by-field map from the Lean hypotheses to the paper proofs is
maintained in the repository at \texttt{docs/core/P9\_PAPER\_PROOF\_MAP.md}.

\section{Normalization and Lean mapping for $b_\rho$ and $H(s)$}\label{app:normalization}
This appendix fixes the normalization of the explicit-formula coefficients and
records the exact Lean correspondence. In the companion preprint we define
\[
  b_\rho := \frac{\rho-1}{\rho}\,G(\rho), \qquad
  G(s) := \pi^{-s/2}\Gamma\!\left(\frac{s}{2}\right)H(s),
\]
where $H(s)$ is a holomorphic factor encoding the ERU kernel normalization and
any auxiliary weights. In the Lean development, $b_\rho$ is an abstract
coefficient attached to a zero, and the explicit-formula identification is
packaged by \texttt{HbExplicitData.h\_b\_rho}. The function
\texttt{explicit\_b\_rho\_expression} is defined as the gamma prefactor
\texttt{gamma\_prefactor}, i.e.\ $\pi^{-\rho/2}\Gamma(\rho/2)$, and any remaining
holomorphic factors are absorbed into the preprint's $H(s)$ and into the abstract
coefficient $b_\rho$ in Lean. Under Assumption~\ref{assm:explicit}, the chosen
normalization matches the Lean predicate \texttt{ERURH.ExplicitBRhoExpression}.

\begin{center}
\begin{tabular}{ll}
\hline
Paper object & Lean name \\
\hline
$\rho=\beta+i\gamma$ & \texttt{ERURH.Alpha.ZeroOfZeta} (fields \texttt{beta}, \texttt{gamma}) \\
$b_\rho$ & \texttt{ERURH.Alpha.b\_rho} \\
$\pi^{-\rho/2}\Gamma(\rho/2)$ & \texttt{ERURH.Alpha.gamma\_prefactor} \\
Explicit identification $b_\rho=\mathrm{explicit\_b}_\rho(\rho)$ &
\texttt{ERURH.Alpha.HbExplicitData.h\_b\_rho} \\
Explicit formula package & \texttt{ERURH.ExplicitBRhoExpression} \\
\hline
\end{tabular}
\end{center}

This normalization is the bridge between the analytic preprint (A/B/C) and the
Lean formalization: the analytic estimates apply to the same coefficients once
the above identification is fixed.

\subsection*{Buchstab kernel factor for the alpha bridge}\label{app:buchstab-kernel}
The alpha bridge data use the multiplicative averaging kernel
$(\mathcal{T}h)(u)=\frac{1}{u}\int_{u/3}^{u/2} h(t)\,dt$. For $h(t)=t^{s-1}$,
the Mellin multiplier is
\[
  m(s)=\frac{2^{-s}-3^{-s}}{s}, \qquad 2^{-s}:=\exp(-s\log 2),
\]
so the kernel contribution to $H(s)$ is $m(s)$. For $\Rea(s)>0$,
Lemma~\ref{lem:buchstab-m-nonzero} gives $m(s)\neq 0$. Thus the only remaining
nonvanishing requirement in Assumption~\ref{assm:explicit} concerns auxiliary
holomorphic weights absorbed into $H(s)$.

\section{Verification of H1--H3 for the ERU kernel}\label{app:h1h3}
This short appendix explains why the hypotheses (H1)–(H3) used in
\texttt{Theorem\_ABC\_preprint.tex} are reasonable for the ERU kernel in the
ERURH setting.

\paragraph{H1 (holomorphic model with gamma factor).}
The explicit-formula package (Assumption~\ref{assm:explicit}) identifies the
coefficients with the completed zeta prefactor
$\pi^{-\rho/2}\Gamma(\rho/2)$ under the chosen normalization. The remaining
kernel-dependent factor is absorbed into $H(s)$ in the preprint, matching the
model $G(s)=\pi^{-s/2}\Gamma(s/2)H(s)$.

\paragraph{H2 (polynomial growth of $H$).}
For the ERU kernel, $H(s)$ is a Mellin transform of a smooth compactly
supported test function and auxiliary weights. Such Mellin transforms are
holomorphic in vertical strips and grow at most polynomially in $|\Im s|$.
This is the standard setting of explicit-formula derivations (see
\cite{Titchmarsh} and the remark on hypothesis verification in the preprint).

\paragraph{H3 (Riemann--von Mangoldt and local zero counts).}
The unit-interval zero count bound is a direct consequence of the classical
Riemann--von Mangoldt formula and is standard; see \cite{Titchmarsh} and
\cite{MontgomeryVaughan}. This is the only zero-distribution input in Theorem B.

\begin{thebibliography}{9}
\bibitem{Titchmarsh}
E. C. Titchmarsh,
\emph{The Theory of the Riemann Zeta-Function},
2nd ed., revised by D. R. Heath-Brown, Oxford University Press, 1986.

\bibitem{MontgomeryVaughan}
H. L. Montgomery and R. C. Vaughan,
\emph{Multiplicative Number Theory I. Classical Theory},
Cambridge University Press, 2007.

\bibitem{Ivic}
A. Ivic,
\emph{The Riemann Zeta-Function: Theory and Applications},
Dover, 2003.

\end{thebibliography}

\end{document}



